\textbf{Phần 1. Trắc nghiệm lựa chọn}
\setcounterpageref{ex}{0}

\begin{ex} %[5D6V35]
Tính nhanh giá trị của biểu thức: $\dfrac{3}{5} \times \dfrac{7}{11} + \dfrac{3}{5} \times \dfrac{4}{11} + \dfrac{2}{5}$.
\choice
{$\dfrac{3}{5}$}
{\True $1$}
{$\dfrac{4}{5}$}
{$2$}
\loigiai{
Ta nhóm hai số hạng đầu:
$$\dfrac{3}{5} \times \left(\dfrac{7}{11} + \dfrac{4}{11}\right) + \dfrac{2}{5} = \dfrac{3}{5} \times 1 + \dfrac{2}{5} = \dfrac{3}{5} + \dfrac{2}{5} = 1.$$
\textbf{Đáp án của câu hỏi là B}
}
\end{ex}

\begin{ex} %[5D6V35]
Tìm $x$, biết: $(x - 4{,}2) : 1{,}8 = 3{,}5$.
\choice
{$9{,}5$}
{\True $10{,}5$}
{$10{,}8$}
{$11{,}2$}
\loigiai{
Ta có:
\begin{align*}
x - 4{,}2 &= 3{,}5 \times 1{,}8\\
x - 4{,}2 &= 6{,}3\\
x &= 6{,}3 + 4{,}2 = 10{,}5
\end{align*}
\textbf{Đáp án của câu hỏi là B}
}
\end{ex}

\begin{ex} %[5D6V35]
Một mảnh đất hình thang có độ dài hai đáy lần lượt là $30\text{ m}$ và $20\text{ m}$, chiều cao $15\text{ m}$. Người ta dùng $20\%$ diện tích mảnh đất để làm nhà ở. Diện tích đất làm nhà ở là
\choice
{$60\text{ m}^2$}
{\True $75\text{ m}^2$}
{$80\text{ m}^2$}
{$150\text{ m}^2$}
\loigiai{
Diện tích mảnh đất là:
$(30 + 20) \times 15 : 2 = 375\text{ (m}^2\text{)}$.\\
Diện tích đất làm nhà ở là:
$375 \times 20\% = 75\text{ (m}^2\text{)}$.\\
\textbf{Đáp án của câu hỏi là B}
}
\end{ex}

\begin{ex} %[5D6V35]
Một người đi xe đạp đi hết quãng đường dài $36\text{ km}$ trong thời gian $2\text{ giờ }15\text{ phút}$. Vận tốc của người đi xe đạp đó là
\choice
{$14\text{ km/h}$}
{$15\text{ km/h}$}
{\True $16\text{ km/h}$}
{$18\text{ km/h}$}
\loigiai{
Đổi: $2\text{ giờ }15\text{ phút} = 2{,}25\text{ giờ}$.\\
Vận tốc người đi xe đạp là:
$36 : 2{,}25 = 16\text{ (km/h)}$.\\
\textbf{Đáp án của câu hỏi là C}
}
\end{ex}

\begin{ex} %[5D6V35]
Một bồn hoa hình tròn có đường kính $3\text{ m}$. Chu vi của bồn hoa đó là
\choice
{$4{,}71\text{ m}$}
{\True $9{,}42\text{ m}$}
{$18{,}84\text{ m}$}
{$7{,}065\text{ m}$}
\loigiai{
Chu vi bồn hoa hình tròn là:
$3 \times 3{,}14 = 9{,}42\text{ (m)}$.\\
\textbf{Đáp án của câu hỏi là B}
}
\end{ex}

\begin{ex} %[5D6V35]
Biết $12\text{ l}$ dầu ăn cân nặng $9{,}6\text{ kg}$. Hỏi một can chứa $18\text{ l}$ dầu ăn đó (không tính khối lượng vỏ can) cân nặng bao nhiêu ki-lô-gam?
\choice
{$13{,}5\text{ kg}$}
{\True $14{,}4\text{ kg}$}
{$15{,}2\text{ kg}$}
{$16\text{ kg}$}
\loigiai{
Một lít dầu ăn cân nặng số ki-lô-gam là:
$9{,}6 : 12 = 0{,}8\text{ (kg)}$.\\
$18\text{ l}$ dầu ăn cân nặng số ki-lô-gam là:
$0{,}8 \times 18 = 14{,}4\text{ (kg)}$.\\
\textbf{Đáp án của câu hỏi là B}
}
\end{ex}

\begin{ex} %[5D6V35]
Một lớp học có $40$ học sinh, trong đó số học sinh nữ chiếm $55\%$. Số học sinh nam của lớp học đó là
\choice
{$22\text{ học sinh}$}
{$20\text{ học sinh}$}
{\True $18\text{ học sinh}$}
{$16\text{ học sinh}$}
\loigiai{
Tỉ số phần trăm học sinh nam là:
$100\% - 55\% = 45\%$.\\
Số học sinh nam là:
$40 \times 45\% = 18\text{ (học sinh)}$.\\
\textbf{Đáp án của câu hỏi là C}
}
\end{ex}

\begin{ex} %[5D6V35]
Tính giá trị của biểu thức: $45{,}8 \times 10 - 258 \times 0{,}1$.
\choice
{$422{,}2$}
{\True $432{,}2$}
{$433{,}2$}
{$442{,}2$}
\loigiai{
Ta có:
$$45{,}8 \times 10 - 258 \times 0{,}1 = 458 - 25{,}8 = 432{,}2.$$
\textbf{Đáp án của câu hỏi là B}
}
\end{ex}

\begin{ex} %[5D6V35]
Một hình tam giác có cạnh đáy dài $16\text{ cm}$. Nếu kéo dài cạnh đáy thêm $4\text{ cm}$ thì diện tích hình tam giác tăng thêm $24\text{ cm}^2$. Diện tích hình tam giác ban đầu là
\choice
{$48\text{ cm}^2$}
{$72\text{ cm}^2$}
{\True $96\text{ cm}^2$}
{$120\text{ cm}^2$}
\loigiai{
Chiều cao của tam giác là:
$24 \times 2 : 4 = 12\text{ (cm)}$.\\
Diện tích tam giác ban đầu là:
$16 \times 12 : 2 = 96\text{ (cm}^2\text{)}$.\\
\textbf{Đáp án của câu hỏi là C}
}
\end{ex}

\begin{ex} %[5D6V35]
Đổi $4\text{ m}^2\text{ }25\text{ cm}^2$ ra đề-xi-mét vuông ta được kết quả là
\choice
{$42{,}5\text{ dm}^2$}
{$425\text{ dm}^2$}
{\True $400{,}25\text{ dm}^2$}
{$40{,}25\text{ dm}^2$}
\loigiai{
Vì $1\text{ m}^2 = 100\text{ dm}^2$ và $1\text{ dm}^2 = 100\text{ cm}^2$ nên:
$4\text{ m}^2 = 400\text{ dm}^2$; $25\text{ cm}^2 = 0{,}25\text{ dm}^2$.\\
Do đó: $4\text{ m}^2\text{ }25\text{ cm}^2 = 400{,}25\text{ dm}^2$.\\
\textbf{Đáp án của câu hỏi là C}
}
\end{ex}

\begin{ex} %[5D6V35]
Bác Ba mua $3\text{ kg}$ gạo tẻ và $2\text{ kg}$ gạo nếp hết tất cả $92\,000$ đồng. Biết giá $1\text{ kg}$ gạo nếp đắt hơn giá $1\text{ kg}$ gạo tẻ là $6\,000$ đồng. Giá tiền $1\text{ kg}$ gạo tẻ là
\choice
{$14\,000\text{ đồng}$}
{\True $16\,000\text{ đồng}$}
{$18\,000\text{ đồng}$}
{$20\,000\text{ đồng}$}
\loigiai{
Nếu thay $2\text{ kg}$ gạo nếp bằng $2\text{ kg}$ gạo tẻ thì số tiền giảm đi là:
$6\,000 \times 2 = 12\,000\text{ (đồng)}$.\\
Giá tiền của $5\text{ kg}$ gạo tẻ là:
$92\,000 - 12\,000 = 80\,000\text{ (đồng)}$.\\
Giá tiền của $1\text{ kg}$ gạo tẻ là:
$80\,000 : 5 = 16\,000\text{ (đồng)}$.\\
\textbf{Đáp án của câu hỏi là B}
}
\end{ex}

\begin{ex} %[5D6V35]
Khi thực hiện phép chia $158 : 3{,}4$, nếu thương lấy đến hai chữ số ở phần thập phân là $46{,}47$ thì số dư của phép chia là
\choice
{$0{,}2$}
{$0{,}02$}
{\True $0{,}002$}
{$2$}
\loigiai{
Ta có: $46{,}47 \times 3{,}4 = 157{,}998$.\\
Số dư là:
$158 - 157{,}998 = 0{,}002$.\\
\textbf{Đáp án của câu hỏi là C}
}
\end{ex}

\textbf{Phần 2. Trắc nghiệm đúng sai}
\setcounterpageref{ex}{0}

\begin{ex} %[5D6V35]
Cho hai biểu thức: $P = 7{,}5 \times 3{,}6 + 7{,}5 \times 6{,}4$ và $Q = 150 : 2{,}5 + 15$.
\choiceTF
{\True Giá trị của biểu thức $P$ bằng $75$}
{\True Giá trị của biểu thức $Q$ bằng $75$}
{Giá trị của biểu thức $P$ lớn hơn giá trị của biểu thức $Q$}
{\True Hiệu $P - Q$ bằng $0$}
\loigiai{
\begin{itemchoice}
\item Đúng. $P = 7{,}5 \times (3{,}6 + 6{,}4) = 7{,}5 \times 10 = 75$.
\item Đúng. $Q = 60 + 15 = 75$.
\item Sai. Vì $P = Q = 75$.
\item Đúng. $75 - 75 = 0$.
\end{itemchoice}
}
\end{ex}

\begin{ex} %[5D6V35]
Một thửa ruộng hình thang có đáy lớn $50\text{ m}$, đáy bé $30\text{ m}$ và chiều cao $25\text{ m}$.
\choiceTF
{\True Diện tích của thửa ruộng là $1\,000\text{ m}^2$}
{\True Diện tích của thửa ruộng tương đương $0{,}1\text{ ha}$}
{\True Nếu cả thửa ruộng thu hoạch được $650\text{ kg}$ thóc thì năng suất đạt $65\text{ kg}$ thóc trên mỗi $100\text{ m}^2$}
{Nếu mở rộng đáy lớn thêm $10\text{ m}$ và giữ nguyên chiều cao thì diện tích thửa ruộng tăng thêm $250\text{ m}^2$}
\loigiai{
\begin{itemchoice}
\item Đúng. Diện tích: $(50 + 30) \times 25 : 2 = 1\,000\text{ (m}^2\text{)}$.
\item Đúng. Vì $1\text{ ha} = 10\,000\text{ m}^2$ nên $1\,000\text{ m}^2 = 0{,}1\text{ ha}$.
\item Đúng. $1\,000\text{ m}^2$ gấp $100\text{ m}^2$ là $10$ lần, nên mỗi $100\text{ m}^2$ thu được: $650 : 10 = 65\text{ (kg)}$.
\item Sai. Diện tích tăng thêm là: $10 \times 25 : 2 = 125\text{ (m}^2\text{)} 
eq 250\text{ m}^2$.
\end{itemchoice}
}
\end{ex}

\textbf{Phần 3. Trắc nghiệm trả lời ngắn}
\setcounterpageref{ex}{0}

\begin{ex} %[5D6V35]
Tính giá trị của biểu thức: $0{,}125 \times 40 \times 2{,}5 \times 8$.
\shortans{100}
\loigiai{
Áp dụng tính chất giao hoán và kết hợp:
$$(0{,}125 \times 8) \times (40 \times 2{,}5) = 1 \times 100 = 100.$$
}
\end{ex}

\begin{ex} %[5D6V35]
Tìm số tự nhiên $x$ thỏa mãn: $15{,}8 < 2{,}5 \times x < 18{,}2$.
\shortans{7}
\loigiai{
Ta có: $15{,}8 : 2{,}5 < x < 18{,}2 : 2{,}5 \implies 6{,}32 < x < 7{,}28$.\\
Vì $x$ là số tự nhiên nên $x = 7$.
}
\end{ex}

\begin{ex} %[5D6V35]
Một bánh xe đạp có đường kính $0{,}5\text{ m}$. Hỏi bánh xe phải quay bao nhiêu vòng để xe đi được quãng đường dài $314\text{ m}$?
\shortans{200}
\loigiai{
Chu vi của bánh xe là:
$0{,}5 \times 3{,}14 = 1{,}57\text{ (m)}$.\\
Số vòng quay của bánh xe là:
$314 : 1{,}57 = 200\text{ (vòng)}$.
}
\end{ex}

\begin{ex} %[5D6V35]
Một mảnh vườn hình chữ nhật có chiều rộng $15\text{ m}$, diện tích bằng diện tích một mảnh đất hình vuông cạnh $30\text{ m}$. Chiều dài của mảnh vườn hình chữ nhật đó là bao nhiêu mét?
\shortans{60}
\loigiai{
Diện tích mảnh đất hình vuông là:
$30 \times 30 = 900\text{ (m}^2\text{)}$.\\
Chiều dài của mảnh vườn hình chữ nhật là:
$900 : 15 = 60\text{ (m)}$.
}
\end{ex}

\begin{ex} %[5D6V35]
Trung bình cộng của hai số bằng $45$. Số thứ nhất bằng $\dfrac{2}{3}$ số thứ hai. Tìm số lớn (số thứ hai).
\shortans{54}
\loigiai{
Tổng của hai số là:
$45 \times 2 = 90$.\\
Số lớn (số thứ hai) là:
$90 : (2 + 3) \times 3 = 54$.
}
\end{ex}

\begin{ex} %[5D6V35]
Một tấm tôn hình tam giác có cạnh đáy dài $30\text{ dm}$ và chiều cao tương ứng bằng $18\text{ dm}$. Diện tích của tấm tôn đó là bao nhiêu đề-xi-mét vuông?
\shortans{270}
\loigiai{
Diện tích của tấm tôn là:
$30 \times 18 : 2 = 270\text{ (dm}^2\text{)}$.
}
\end{ex}
